%% ch07-incident-response.tex — JA4proxy Operator's User Guide
%% Chapter 7: Incident Response

\chapter{Incident Response}
\label{ch:incident}
\index{incident response}

\section{Accessing the Dashboard Remotely}
\index{SSH tunnel}
\index{dashboard!remote access}

\newthought{The management dashboard is bound to loopback only} (\code{127.0.0.1}).
In a typical deployment the proxy runs in a datacentre DMZ host and you are on a
laptop at home or over a corporate VPN. You cannot open \code{localhost:8090} in a
browser — you need an SSH port-forward first.

\subsection{Setting Up the SSH Tunnel}

If your organisation routes SSH through a bastion host (common in regulated
environments):

\begin{bashcode}
ssh -J YOU@bastion.mgmt.corp.example.com YOU@dmz-web01 \
    -L 8090:127.0.0.1:8090 \
    -L 8444:127.0.0.1:8444 \
    -N &
\end{bashcode}

If your VPN gives you direct SSH access to DMZ hosts:

\begin{bashcode}
ssh YOU@dmz-web01.corp.example.com \
    -L 8090:127.0.0.1:8090 \
    -L 8444:127.0.0.1:8444 \
    -N &
\end{bashcode}

Then open in your browser:

\begin{itemize}
  \item \code{https://127.0.0.1:8444} — HTTPS (Caddy sidecar, self-signed cert;
    click through the certificate warning — this is expected)
  \item \code{http://127.0.0.1:8090} — plain HTTP fallback if the HTTPS sidecar
    is not running
\end{itemize}

\begin{tipbox}[Corporate proxy and localhost]
Some corporate web proxies intercept \code{localhost} URLs. Use
\code{127.0.0.1} (the numeric address) in every URL — it bypasses the proxy.
\end{tipbox}

For two-datacentre setups and full troubleshooting guidance, see
\code{docs/runbooks/dashboard\_access.md}.

\section{Emergency Deploy (JA4proxy Not Yet Running)}
\index{emergency deploy}

If \ja\ is not yet deployed on the host, get it running in under 5 minutes:

\begin{bashcode}
# SSH into the DMZ host
ssh -J YOU@bastion YOU@dmz-web01

cd /opt && git clone https://github.com/seanpor/JA4proxy.git && cd JA4proxy
export BACKEND_HOST=your-actual-backend.internal.example.com

# Pull pre-built images — no compiler or build tools needed on the host
docker compose pull && docker compose up -d
\end{bashcode}

Pre-built images are published to GHCR on every merge to main:
\begin{itemize}
  \item \code{ghcr.io/seanpor/ja4proxy-go:main}
  \item \code{ghcr.io/seanpor/ja4proxy-management:main}
\end{itemize}

Verify the proxy is up before proceeding:
\begin{bashcode}
docker compose ps          # all containers should be "Up"
curl -kv https://127.0.0.1:8443/   # should return backend response
\end{bashcode}

\section{Inserting JA4proxy Into the Live Traffic Path}
\index{traffic insertion}
\index{iptables}

If \ja\ is running but traffic is not yet flowing through it, insert it with a
single iptables rule:

\begin{bashcode}
make traffic-on
# Runs: sudo iptables -t nat -I PREROUTING -p tcp --dport 443 \
#         -j REDIRECT --to-port 8443
\end{bashcode}

This redirects all inbound port 443 connections to JA4proxy on port 8443.
It takes effect immediately with no service restart. Instant rollback:

\begin{bashcode}
make traffic-off
# Restores the original path in under 1 second.
\end{bashcode}

\begin{notebox}[iptables availability]
\code{make traffic-on} requires \code{sudo iptables} on the DMZ host. If
iptables is not available (some hardened or cloud hosts restrict it), use a
HAProxy backend swap or DNS cutover instead — see
\code{docs/runbooks/dashboard\_access.md\#step-5}.
\end{notebox}

\newthought{An active attack looks different depending on its type.}
A fingerprint-consistent botnet shows up immediately in the Grafana fingerprint
panel. A slow scan with randomised tool configurations might only show up in the
analytics node's campaign detection. A sudden volumetric spike shows up in the
connection rate panel within seconds.

This chapter gives you a structured approach to detecting, responding to, and
recovering from incidents involving \ja.

\section{Detecting an Active Attack}
\index{attack detection}

\subsection{What to Look For in Grafana}

Open the \emph{JA4proxy Overview} dashboard and look for these indicators:

\begin{description}[leftmargin=4.5cm]
  \item[Connection rate spike] A sudden increase in connections/second, especially
    if the action distribution shows a high BLOCKED fraction. This is a good sign
    — the proxy is detecting and handling the attack.

  \item[Risk score distribution shift] The histogram shifts right (higher scores).
    This indicates more traffic scoring as suspicious.

  \item[New dominant fingerprint] A single fingerprint appears at the top of the
    Top Fingerprints table with high connection counts and high average scores.

  \item[Geographic concentration] The map shows a sudden spike from a region
    you do not normally serve.

  \item[Beaconing spikes] The analytics node detects a cluster of IPs beaconing
    at similar intervals — characteristic of C2 callback traffic.
\end{description}

\subsection{Using the Score Drift Alert}

The \code{ScoreDrift} Alertmanager alert fires when the mean risk score across all
connections drifts more than 20 points from the 24-hour baseline. This catches
attacks that do not produce obvious connection rate spikes — such as slow scanners
or distributed credential-stuffing campaigns spread across many source IPs.

\subsection{Checking the Analytics Node}

The analytics node performs cross-instance aggregation and campaign detection.
Its findings are pushed to Redis and read by the scorer. Check its logs for
detected patterns:

\begin{bashcode}
docker compose logs analytics | grep -E "campaign|slow.scan|score.drift"
\end{bashcode}

\section{Emergency Blocking Procedures}
\index{emergency blocking}
\index{blacklist!emergency}

\subsection{Blocking a JA4 Fingerprint Immediately}
\index{ja4-admin.sh}

When you identify a malicious fingerprint and need to block it immediately, use
the admin CLI — it takes effect for new connections within seconds:

\begin{bashcode}
./scripts/ja4-admin.sh block-ja4 t13d0912_a1b2c3d4e5f6_987654321abc
\end{bashcode}

Verify the block is active:
\begin{bashcode}
docker compose logs -f ja4proxy | grep "JA4 blacklisted"
\end{bashcode}

You can also block directly from the dashboard: open the \emph{Top Fingerprints}
panel and click the \textbf{Block} button next to the fingerprint.

\begin{tipbox}[Safety guarantee — what to say on the incident call]
\emph{``Real browsers using Chrome, Firefox, or Safari cannot be blocked by this
proxy. Browser TLS fingerprints always include h2/h1 ALPN and are bypassed before
scoring. The fingerprint we are blocking is a Python library signature with no
browser ALPN. Legitimate users on the form are unaffected.''}
\end{tipbox}

\begin{notebox}[Persistent blocking across restarts]
Blocking via the CLI or dashboard writes to Redis and takes effect immediately,
but is lost on a Redis restart. For a permanent entry that survives restarts,
also add the fingerprint to \code{config/proxy.yml}:
\begin{yamlcode}
security:
  blacklist:
    - t13d0912_a1b2c3d4e5f6_987654321abc  # Sliver C2 — 2026-04-04 incident
\end{yamlcode}
Then reload config:
\begin{bashcode}
docker compose kill -s HUP ja4proxy
\end{bashcode}
\end{notebox}

\subsection{Blocking a Specific IP Address}

To immediately ban an IP for the standard 5-minute TTL:

\begin{bashcode}
docker compose exec redis redis-cli SETEX "ban:185.220.0.1" 300 "manual"
\end{bashcode}

For a longer ban (e.g., 1 hour = 3600 seconds):

\begin{bashcode}
docker compose exec redis redis-cli SETEX "ban:185.220.0.1" 3600 "manual-1h"
\end{bashcode}

\begin{notebox}[Ban TTLs and self-healing]
There are no permanent bans in \ja. The maximum practical ban duration is
whatever you set manually in Redis — but be careful with long ban durations.
Shared IP infrastructure (NAT, corporate proxies, CDN exit nodes) means long
bans can affect many legitimate users sharing an IP.
\end{notebox}

\subsection{Blocking a CIDR Range}

\ja\ does not currently support native CIDR banning via Redis (CIDR matching
is in-process via a trie, not Redis). For immediate CIDR blocking, use HAProxy's
ACL mechanism or your network firewall.

For example, to block \code{185.220.0.0/24} at the HAProxy level:
\begin{lstlisting}[language=bash, numbers=none]
# Add to haproxy.cfg in the frontend section:
acl blocked_range src 185.220.0.0/24
tcp-request connection reject if blocked_range
\end{lstlisting}

Then reload HAProxy:
\begin{bashcode}
docker compose kill -s HUP haproxy
\end{bashcode}

\subsection{Emergency Country Block}

To block an entire country immediately via config:

\begin{yamlcode}
country_blacklist_enabled: true
country_blacklist:
  - RU   # Added 2026-04-04 — active attack campaign
\end{yamlcode}

Reload config. The block takes effect for new connections immediately.

\section{Investigating a Suspicious Fingerprint}
\index{fingerprint investigation}

\subsection{Step 1: Identify the Fingerprint}

From Grafana's Top Fingerprints panel, or from logs:

\begin{bashcode}
docker compose logs ja4proxy | grep "Score: [7-9][0-9]\|Score: 100" | \
  grep -oP "JA4: \K[^ |]+" | sort | uniq -c | sort -rn | head -20
\end{bashcode}

\subsection{Step 2: Look It Up}

Search the JA4 community database: \url{https://ja4db.com}

The database contains community-identified labels for thousands of fingerprints,
including known C2 frameworks, scanners, and security tools.

\subsection{Step 3: Query Redis for Activity}

Check the rate-limit history for this fingerprint:

\begin{bashcode}
docker compose exec redis redis-cli \
  ZCARD "rl:ja4:t13d0912_a1b2c3d4e5f6_987654321abc"
\end{bashcode}

Check beaconing records:

\begin{bashcode}
docker compose exec redis redis-cli --scan \
  --pattern "beacon:*:t13d0912_a1b2c3d4e5f6_987654321abc"
\end{bashcode}

\subsection{Step 4: Assess Risk}

Ask:
\begin{itemize}
  \item Is this fingerprint in the JA4 database as a known malicious tool?
  \item What is its average risk score? (Grafana fingerprint panel)
  \item How many source IPs are using it? (Redis PFCOUNT on the HLL key)
  \item Is it beaconing at regular intervals?
  \item Is it associated with known-bad ASNs or countries?
\end{itemize}

If the answer to the first question is yes and the fingerprint is not shared with
any legitimate tool, add it to the blacklist. If you are unsure, add it to the
watchlist and continue monitoring for 24--48 hours.

\section{Escalating to a Permanent Country Block}
\index{country blocking}

Country blocking is a significant decision — it will affect all users from that
country, including legitimate ones. Follow this procedure:

\begin{enumerate}
  \item Verify the attack is predominantly from the target country (Grafana
    geographic distribution panel).
  \item Check how much legitimate traffic you receive from that country
    (filter Grafana for the country code and check the action distribution).
  \item If legitimate traffic from that country is $<$1\% of total, the block
    has acceptable false-positive impact.
  \item Add to the blacklist in \code{config/proxy.yml}. Note: if
    \configkey{country\_whitelist\_enabled} is true, add the excluded countries
    to the whitelist instead.
  \item Reload config and monitor the block rate for 30 minutes.
  \item Document the decision: which country, why, what the evidence was, who
    approved it.
\end{enumerate}

\section{Recovering From a False Positive}
\index{false positive!recovery}
\index{false positive}

A false positive occurs when a legitimate connection is blocked or banned. The
most common causes are:
\begin{itemize}
  \item A legitimate tool uses a fingerprint that also appears in the blacklist
  \item A legitimate IP has been auto-banned due to high request volume
  \item A country block is affecting legitimate users
\end{itemize}

\subsection{Immediate Recovery}

If a user is currently blocked:

\begin{bashcode}
# Clear their ban immediately
docker compose exec redis redis-cli DEL "ban:192.0.2.1"
\end{bashcode}

If a fingerprint was incorrectly blacklisted:

\begin{enumerate}
  \item Remove the fingerprint from \code{security.blacklist} in
    \code{config/proxy.yml}
  \item Reload config
  \item If the fingerprint was also added to the Redis blacklist directly,
    remove it there too:
\begin{bashcode}
docker compose exec redis redis-cli SREM ja4:blacklist \
  "t13d1113h2_8daaf6152771_b0da82dd1658"
\end{bashcode}
\end{enumerate}

\subsection{Root Cause and Prevention}

After clearing the false positive:

\begin{enumerate}
  \item Identify what triggered it (which bypass condition or scoring signal)
  \item If a specific fingerprint was too broadly blacklisted, consider whether
    the problematic tool shares a fingerprint with legitimate software
  \item If the rate limit was too aggressive, adjust the threshold in
    \code{config/proxy.yml} and reload
  \item Add the legitimate fingerprint to the whitelist to prevent recurrence
  \item Document the incident and the configuration change
\end{enumerate}

\section{Post-Incident Review}
\index{post-incident review}

After resolving any significant incident:

\begin{enumerate}
  \item Export Grafana panel data for the incident time window (use the panel's
    share/export function)
  \item Pull the log segment from Loki covering the incident
  \item Record:
    \begin{itemize}
      \item When it started and when it was resolved
      \item The attack type (fingerprint-based, volumetric, geographic)
      \item Which signals detected it first
      \item What actions were taken
      \item Whether any false positives occurred
    \end{itemize}
  \item Update the fingerprint blacklist, labels, or rate-limit thresholds based
    on what you learned
  \item Consider whether a new alert rule would have caught this faster
\end{enumerate}

\begin{tipbox}[Preserving incident data]
Grafana's snapshot feature captures the current state of the dashboard for a
specific time range. Save a snapshot for significant incidents before the data
rolls out of the retention window:
\begin{lstlisting}[language=bash, numbers=none]
Dashboard menu → Share → Snapshot → Publish
\end{lstlisting}
\end{tipbox}
